\begin{figure}[p]\centering\begin{tikzpicture}[gclplate]\node[anchor=west,font=\sffamily\bfseries\Large,gcllabel] at (-6,3.8){CLASSICAL PROOF AS CONTROL};\draw[GCLWarm,line width=1pt](-6,3.45)--(6,3.45);\node[gcllabel,draw](l) at (-3,1.2){classical proof step};\node[gcllabel,draw](c) at (3,1.2){continuation manipulation};\draw[gclwarmarrow](l)--node[above,gcllabel]{interpretation}(c);\end{tikzpicture}\caption{\textbf{Plate 26. Classical proof as control.} Selected classical principles receive computational readings through continuation structure. Scope: one typed interpretation among several, not universal identity.}\label{plate:classical-control}\end{figure}
